% NỘI DUNG FILE: contents/Task1.tex

Trong hệ thống nhúng thời gian thực (Real-Time Embedded System), việc đảm bảo tính nhất quán của dữ liệu giữa các tiến trình (tasks) đồng thời tránh các xung đột tài nguyên phần cứng là một yêu cầu tối quan trọng. Task 1 tập trung vào việc xây dựng cơ chế thu thập dữ liệu môi trường (nhiệt độ, độ ẩm) và đồng bộ hóa thông tin này để điều khiển các thiết bị cảnh báo trực quan (LED và màn hình LCD) thông qua hệ điều hành thời gian thực FreeRTOS.

Kiến trúc phần mềm của hệ thống cảnh báo được chia thành hai luồng tiến trình độc lập hoạt động song song theo mô hình Nhà sản xuất - Người tiêu thụ (Producer - Consumer), được đồng bộ thông qua các cơ chế Giao tiếp liên tiến trình (Inter-Process Communication - IPC) bao gồm Hàng đợi (Queue) và Cờ hiệu (Semaphore).

\subsubsection{Cơ chế thu thập dữ liệu và cảnh báo qua LCD (Producer Task)}
Tiến trình \texttt{temp\_humi\_monitor} đóng vai trò là luồng thu thập dữ liệu chính (Producer). Tiến trình này giao tiếp với cảm biến DHT20 thông qua chuẩn giao tiếp I2C để đọc các giá trị nhiệt độ và độ ẩm của môi trường. 

Để đảm bảo hệ thống không bị "đóng băng" (blocking) trong quá trình chờ đọc dữ liệu, thay vì sử dụng hàm \texttt{delay()} truyền thống, hệ thống áp dụng cơ chế đếm tick của FreeRTOS (\texttt{xTaskGetTickCount()}). Dữ liệu được lấy mẫu với chu kỳ $T_{sample} = 5000\text{ms}$. 

Ngay sau khi quá trình đọc hoàn tất, dữ liệu hợp lệ sẽ được đóng gói vào cấu trúc \texttt{SensorData\_t} và đẩy vào hàng đợi \texttt{sensorQueue} thông qua API \texttt{xQueueOverwrite()}. Việc sử dụng \texttt{Overwrite} thay vì \texttt{Send} tiêu chuẩn đảm bảo rằng hàng đợi luôn chứa bản cập nhật dữ liệu mới nhất, tránh độ trễ (latency) trong việc xử lý thông tin môi trường.

Bên cạnh đó, tiến trình này tích hợp một Máy trạng thái hữu hạn (Finite State Machine - FSM) để quản lý giao diện hiển thị LCD:
\begin{itemize}
    \item \textbf{Phân loại trạng thái (State Evaluation):} Dữ liệu được đưa qua hàm \texttt{evaluate\_lcd\_state} để đối chiếu với các ngưỡng tĩnh đã được định nghĩa trước (ví dụ: \texttt{TEMP\_HOT}, \texttt{TEMP\_COLD}, \texttt{HUMIDITY\_VERY\_HUMID}). FSM xuất ra ba trạng thái hiển thị chính: \texttt{NORMAL}, \texttt{WARNING}, và \texttt{CRITICAL}.
    \item \textbf{Xử lý ngoại lệ hiển thị:} Trong điều kiện bình thường, LCD hiển thị giá trị thời gian thực. Khi rơi vào trạng thái bất thường (Abnormal), hệ thống tự động chèn thêm chuỗi cảnh báo "SOS" vào màn hình hiển thị.
    \item \textbf{Điều khiển đèn nền động:} Để thu hút sự chú ý của người dùng trong trạng thái \texttt{CRITICAL}, một FSM phụ trợ (\texttt{handle\_lcd\_backlight\_fsm}) được kích hoạt. FSM này điều khiển nhấp nháy đèn nền LCD với chu kỳ 500ms mà không làm gián đoạn chu kỳ lấy mẫu 5 giây của cảm biến, thể hiện tính ưu việt của lập trình phi đồng bộ.
\end{itemize}

\subsubsection{Cơ chế chỉ thị trạng thái qua LED (Consumer Task)}
Tiến trình \texttt{led\_blinky} hoạt động như một thực thể tiêu thụ dữ liệu (Consumer). Nó liên tục giám sát hàng đợi \texttt{sensorQueue} thông qua hàm \texttt{xQueuePeek()} để sao chép dữ liệu nhiệt độ mới nhất mà không làm rỗng hàng đợi, cho phép các tiến trình khác trong hệ thống (như tiến trình đẩy dữ liệu lên IoT Cloud) tiếp tục sử dụng lượng dữ liệu này.

Tần số nhấp nháy của LED cảnh báo không cố định mà là một hàm phụ thuộc vào mức độ nghiêm trọng của nhiệt độ môi trường, tuân theo luật điều khiển sau:
\begin{table}[H]
    \centering
    \caption{Phân rã chu kỳ nhấp nháy LED theo mức cảnh báo nhiệt độ}
    \label{tab:led_intervals}
    \begin{tabular}{|l|l|c|}
    \hline
    \textbf{Trạng thái FSM} & \textbf{Điều kiện kích hoạt} & \textbf{Chu kỳ chớp (ms)} \\ \hline
    \texttt{LED\_STATE\_TEMP\_HOT}  & $T \geq T_{HOT}$ & 100 \\ \hline
    \texttt{LED\_STATE\_TEMP\_WARM} & $T_{WARM} \leq T < T_{HOT}$ & 300 \\ \hline
    \texttt{LED\_STATE\_NORMAL\_TEMP} & $T_{COOL} < T < T_{WARM}$ & 1000 \\ \hline
    \texttt{LED\_STATE\_TEMP\_COOL} & $T_{COLD} < T \leq T_{COOL}$ & 2000 \\ \hline
    \texttt{LED\_STATE\_TEMP\_COLD} & $T \leq T_{COLD}$ & 3000 \\ \hline
    \end{tabular}
\end{table}

Động lực học của thiết kế này tuân theo nguyên lý Ergonomics trong cảnh báo công nghiệp: tần số chớp (frequency) tỷ lệ thuận với mức độ khẩn cấp, giúp người vận hành nhận biết tức thời trạng thái nguy hiểm của hệ thống.

\subsubsection{Tránh tranh chấp tài nguyên bằng Semaphore}
Một đặc điểm kỹ thuật nổi bật trong kiến trúc của hệ thống là việc giải quyết bài toán tranh chấp điều khiển chân GPIO của LED. Vì LED có thể được điều khiển tự động bởi nhiệt độ (tiến trình \texttt{led\_blinky}), nhưng cũng có thể bị can thiệp ghi đè bởi người dùng thông qua giao diện Web/WebSocket (tiến trình \texttt{task\_handler}).

Để giải quyết vấn đề này, hệ thống áp dụng cơ chế khóa tài nguyên Mutex/Semaphore (\texttt{xSemaphoreLedControl}). 
\begin{itemize}
    \item \textbf{Cơ chế cấp phép:} Trước khi thực hiện chuyển đổi trạng thái mức logic của chân \texttt{LED\_GPIO}, tiến trình bắt buộc phải gọi hàm \texttt{xSemaphoreTake(..., 0)}. Nếu lấy được cờ hiệu (thành công trả về \texttt{pdTRUE}), tiến trình tiến hành chớp tắt LED bình thường và lập tức hoàn trả cờ hiệu bằng \texttt{xSemaphoreGive()}.
    \item \textbf{Cơ chế ghi đè (Override):} Nếu một tác vụ ngoại vi đang giữ cờ hiệu (nghĩa là đang có lệnh bắt buộc tắt LED từ Server), hàm \texttt{Take} sẽ thất bại. Luồng điều khiển sẽ nhảy vào nhánh \texttt{else}, tự động kéo chân GPIO xuống mức \texttt{LOW} và cập nhật lại bộ đếm thời gian.
\end{itemize}

Cách tiếp cận này không chỉ bảo vệ tài nguyên phần cứng khỏi các trạng thái race-condition (điều kiện tương tranh) giữa các vi luồng, mà còn đảm bảo hệ thống có thể linh hoạt chuyển đổi giữa chế độ hoạt động Tự động (Auto/Sensor-driven) và Thủ công (Manual/User-driven) một cách an toàn và mượt mà.